\chapter{Warm Starts and Families of Problems}
\label{ch:warmstart}

Solvers are usually studied one problem at a time. Applications almost never present
them that way. An efficient frontier is a family of QPs differing only in a return
tilt. A rolling rebalance is a family differing only in the data. A scenario grid, a
cross-validation sweep, a sensitivity study, a resampled cloud: each is dozens or
hundreds of instances that are nearly the same.

Solved independently, each instance rediscovers an active set it almost always
already had. This chapter is about not doing that. It has three parts: what can be
\emph{reused} exactly, what must be \emph{repaired}, and why neither can change the
answer.

\section{What does not depend on the linear term}

Start with an observation about the invariants of Chapter~\ref{ch:walking} that is
easy to read past.
\begin{equation*}
  J\T G J = I_n, \qquad J\T A = \begin{bmatrix} R \\ 0 \end{bmatrix}.
\end{equation*}
\emph{Neither factor depends on $a$.} $J$ depends on $G$ alone, through
$J\T GJ = I$, and on the working set through the orthogonal transformations
accumulated into it; $R$ depends on $G$ and the working set. The linear term enters
only through $x_u = G^{-1}a$ and the right-hand sides.

So across a family of programs in which $G$, $C$, $b$ are fixed and only $a$ varies,
\emph{both factors are reusable verbatim}, and the entire solution can be recovered
from them.

\begin{proposition}[Recovery from cached factors]
\label{prop:recover}
\index{warm start!recovery}
Let $(J,R)$ satisfy the invariants for a working set $\Active$ of size $k$ with
normals $A$. For a new linear term $a$, set
\begin{equation}
  x_u = J\,(J\T a),
  \qquad
  \rho = b_\Active - A\T x_u,
  \qquad
  R\T y = \rho,
  \qquad
  x = x_u + J_1 y,
  \qquad
  R\,u = y .
  \label{eq:recover}
\end{equation}
Then $(x,u)$ is the solution and multiplier pair of the working-set
subproblem~\eqref{eq:sub}.
\end{proposition}

\begin{proof}
$x_u = JJ\T a = G^{-1}a$ by the first invariant. By
Proposition~\ref{prop:blocks}(i), $u = R^{-1}y = (R\T R)^{-1}\rho = (A\T G^{-1}A)^{-1}\rho$,
the multiplier of~\eqref{eq:subsol}. For the iterate, $Ru = y$ gives
$J_1 y = J_1 R u = G^{-1}Au$ by Proposition~\ref{prop:blocks}(iii), so
$x = x_u + G^{-1}Au$, again as in~\eqref{eq:subsol}.
\end{proof}

\subsection{What the recovery costs, and a correction worth making}

The cost is $O(n^2)$, \emph{all of it in $x_u = J(J\T a)$}: two matrix--vector products
with a dense $n\times n$ matrix, $J$ having lost its triangularity at the first update.
Everything after that is $O(nk + k^2)$. Re-deriving the factors for the same working
set instead costs $k$ insertions at $O(n^2)$ each, so the saving is a factor of order
$k$---largest exactly where it matters, since a long-only optimum is a vertex at which
most bounds bind.

The exponent is easy to get wrong, and the reason is instructive.

\begin{remark}[The natural experiment cannot see the exponent]
\label{rem:exponent}
On a box-constrained family the active set grows with $n$, so $k = \Theta(n)$ and the
predictions $O(nk)$ and $O(n^2)$ describe the \emph{same curve}. Either reading fits
the data, and a paper can report ``linear in $nk$'' without anything looking wrong.

Separating them requires holding $k$ \emph{fixed} while $n$ grows. Doing so, the
measured local exponent $d\log t/d\log n$ runs from well below one at small $n$ to
close to two at large $n$---and neither end is $1$. The flat small-$n$ end is not a
different asymptotic; it is an \emph{overhead floor}, where a solve costs what its
dozen array operations cost to dispatch, and it ends where the $n^2$ arithmetic
overtakes dispatch.

The general lesson is one Chapter~\ref{ch:measurement} develops: \emph{a benchmark
family in which two candidate exponents coincide cannot distinguish them}, and
designing the family that breaks the tie is part of the work.
\end{remark}

There is also an asymmetry in cost between recovering and \emph{checking}. The recovery
is $O(n^2)$, but the certificate must look at every constraint, so it is $O(nm)$ for a
dense $C$. On a dense problem with many constraints, \emph{verification and not
recovery dominates a successful reuse}. That is a slightly deflating fact and a useful
one: it tells you where to spend optimisation effort.

\section{Repair rather than restart}

When the certificate fails, the cached working set is stale---but it remains a far
better starting point than the unconstrained minimum. What blocks resuming from it is
the dual method's invariant~\eqref{eq:invariant}: the iteration requires dual
feasibility, and a stale set typically has one or more negative multipliers. Those are
precisely the constraints that no longer belong.

\begin{proposition}[Repair terminates in a resumable state]
\label{prop:repair}
\index{warm start!repair}
Iterate the following from the cached set: recover $(x,u)$ by~\eqref{eq:recover}; if
every inequality multiplier is non-negative, stop; otherwise drop one constraint with a
negative multiplier, downdate $(J,R)$, and repeat. The loop stops after at most $k$
passes, in a state satisfying~\eqref{eq:invariant}.
\end{proposition}

\begin{proof}
Each pass that does not stop shrinks the working set by one, so there are at most $k$;
with an empty set the sign condition is vacuous, which is the cold start of
Remark~\ref{rem:cold}. On exit $(x,u)$ solves the subproblem by
Proposition~\ref{prop:recover} and the sign conditions hold, which
is~\eqref{eq:invariant}; full column rank is inherited, deleting columns preserving it.
\end{proof}

From such a state \emph{the iteration cannot tell a resumed solve from a cold one}, the
invariant being all it reads, so the resumed walk needs only to drive the remaining
primal infeasibility to zero. In the worst case everything is dropped and the resumed
walk is a cold solve, so repair is never worse than restarting up to the cost of the
drops themselves.

\section{Warm starts in the other two strategies}

The same idea appears in each of the book's strategies, with a different mechanism.

\paragraph{Block pivoting (Chapter~\ref{ch:guessing}).} The warm start is the previous
problem's \emph{free set} as the initial guess. Since the repair rule needs no dual
feasibility to begin---it reads signs and toggles---no repair loop is required at all;
a stale set simply costs an extra exchange or two.

\paragraph{Krylov inner solves (Chapter~\ref{ch:krylov}).} Two distinct warm starts are
available and they are worth distinguishing, because only one is available to a direct
solver. The \emph{outer} warm start is the free set, which cuts the number of outer
steps. The \emph{inner} warm start is the previous solution as CG's initial iterate,
which cuts inner iterations---and by Proposition~\ref{prop:warm} the inner count then
scales with $\log\norm{\delta}$, the logarithm of the \emph{parameter step}, so a finer
grid is cheaper per point. A direct inner solver profits only from the first, having no
analogue of an initial guess. That is where warm-started parametric sweeps obtain their
largest speed-ups.

\paragraph{The homotopy (Chapter~\ref{ch:parametric}).} Here warm starting is not an
optimisation but the algorithm itself: every step is a warm start from the previous
breakpoint. Which is one way of saying that the parametric method is what you get when
you take warm starting to its limit and let the parameter vary continuously.

\begin{remark}[Why active-set methods warm-start and interior-point methods do not]
An active-set method terminates \emph{at} a face of the polyhedron, with the face
recorded as a discrete object---a set of indices. That object is exactly what transfers
to a neighbouring problem. An interior-point method approaches a face asymptotically
from the interior and holds no such object; restarting it near a previous solution puts
the iterate close to the boundary, where the barrier is badly behaved, and the standard
remedy is to back off toward the centre---discarding most of the information one hoped
to reuse. This asymmetry, not raw speed at a single solve, is the main reason
parametric and rolling applications are built on active-set methods.
\end{remark}

\section{Structural continuity: why the active set barely moves}

Warm starting works because consecutive solutions share almost all of their active set.
For a parametric family that is not a hope but a theorem, and it is
Chapter~\ref{ch:parametric}'s.

Along a frontier the solution is piecewise linear in $\lambda$ with generically
\emph{single-index} breakpoint changes. So adjacent solves on a fine grid share all but
a few active constraints, and a warm solver typically terminates in one to two outer
steps. \emph{This is a structural property of parametric quadratic programming, not an
artefact of the data.}

When the family varies in the data rather than in a parameter---a rolling window, a
resampled draw---no such theorem is available, and the empirical picture takes over: the
warm-started outer count tracks the support drift $|\Free \triangle \Free'|$ and
collapses to one whenever the drift is zero. That is the honest statement, and it is
worth keeping the two cases separate rather than claiming the theorem covers both.

\section{Why none of this can change the answer}

Neither reuse nor repair introduces an approximation, and the argument is the
certificate principle of Chapter~\ref{ch:certificate} once more.

Reuse returns a point \emph{only when it carries a certificate}: the working-set
inequalities' multipliers must be non-negative and no constraint outside the working set
may be violated, which by Proposition~\ref{prop:sufficient} is a proof. Repair merely
chooses a different starting state for an iteration whose output is determined
independently of where it started.

What \emph{does} change is the reported iteration count, which counts working-set edits
actually performed and is therefore not comparable with a cold solve's. A benchmark
table mixing warm and cold iteration counts in one column is meaningless, and
Chapter~\ref{ch:measurement} says what to report instead.

\begin{example}[A frontier sweep, end to end]
Fifty points on a long-only efficient frontier for a few hundred assets. Cold, each
point is an independent solve: fifty Cholesky factorisations, fifty walks of tens of
steps. Warm, the first point is a cold solve and each subsequent point recovers from the
cached factors, verifies, and---on the minority of points where the active set has
moved---repairs and resumes for one or two steps. The measured advantage of the warm
sweep over an assemble-and-factor baseline is roughly fivefold, and roughly fifteenfold
against a cold baseline, growing with $n$. Against a general-purpose modelling layer
called once per point, the end-to-end gap is three orders of magnitude---but that figure
mostly measures the modelling layer, and Chapter~\ref{ch:measurement} explains why one
should say so.
\end{example}

\begin{notes}
Propositions~\ref{prop:recover} and~\ref{prop:repair}, and the corrected cost exponent
of Remark~\ref{rem:exponent}, are from the Goldfarb--Idnani note \cite{schmelzer2026quadprog}; the warm-start analysis for
the block-pivoting loop with Krylov inner solves is \cite{schmelzer2026nncg}
(Proposition~\ref{prop:warm}); the frontier-sweep measurements are from
\cite{schmelzer2026rmt} and \cite{schmelzer2026matrixfree}. Parametric warm starting of
active-set QP solvers is treated generally in \cite{nocedal2006}; combining it with a
matrix-free Krylov inner solver on a changing active set is more recent.
\end{notes}

\begin{exercises}

\begin{exercise}
Verify Proposition~\ref{prop:recover} by substituting~\eqref{eq:recover}
into~\eqref{eq:subsol} directly, without using Proposition~\ref{prop:blocks}.
\end{exercise}

\begin{exercise}
Show that the recovery~\eqref{eq:recover} is $O(n^2)$ and identify precisely which two
operations carry the $n^2$. Why does $J$ lose its triangularity after the first update?
\end{exercise}

\begin{exercise}
Prove that the repair loop of Proposition~\ref{prop:repair} preserves full column rank
of the working-set normals, and explain why the corresponding claim would fail for an
\emph{addition} loop.
\end{exercise}

\begin{exercise}
Suppose the cached working set is optimal for the new problem. How many operations does
a successful reuse cost, in terms of $n$, $m$, $k$? Which term dominates for
$m \approx n$, and for $m \approx n^2$?
\end{exercise}

\begin{exercise}
Give a family of problems in which the cached active set is \emph{always} stale, so that
repair drops everything and the warm start is pure overhead. What does this say about
when to attempt reuse?
\end{exercise}

\begin{exercise}
Explain why warm-started and cold iteration counts are not comparable, and propose two
quantities that \emph{are} comparable across the two regimes.
\end{exercise}

\begin{exercise}[Implementation]
Take your solver from Chapter~\ref{ch:walking} and add factor caching plus the recovery
and repair of this chapter. Sweep the linear term over $50$ values and record, per
point: whether reuse succeeded, how many constraints repair dropped, and how many steps
the resumed walk took. Plot all three against the sweep index.
\end{exercise}

\begin{exercise}[Implementation]
Reproduce Remark~\ref{rem:exponent}. Time a warm recovery on a box family where $k$
grows with $n$, fit a power law, and record the exponent. Then hold $k$ fixed at a small
value while $n$ grows and refit. Report both exponents and the $n$ at which the second
family leaves its overhead floor.
\end{exercise}

\end{exercises}
